% !TEX root = journal_v2_0421052099.tex
%
% Packages and styles for the two-column journal version of the thesis.
%
% The class is the standard article class in two columns. Its type sizes copy
% two published papers the author chose as models. Both are kept in
% thesis_templates/references/.
%   - "A survey of graph meets large language" (IJCAI 2024) gives the title,
%     author, and affiliation styles, and most sizes. Body, abstract, and
%     reference text are 10 pt. Captions and footnotes are 9 pt. Section heads
%     are 12 pt bold and subsection heads 11 pt bold.
%   - "KBQA-o1" (ICML 2025) gives the 1 pt rules above and below the title,
%     with its spacing of 0.25 in, 0.22 in, and 0.3 in. Its tables are set at
%     8 pt or 9 pt, and so are these.
% Both papers are on US Letter. This one keeps A4, with the survey's 7 in text
% width and 0.25 in column gap.

\usepackage[a4paper, textwidth=7in, columnsep=0.25in, top=1in, bottom=1in,
  heightrounded]{geometry}

% Times text with Computer Modern maths, as in both model papers and the book.
% The book marks identifiers in sans and scales Helvetica to 92% so that it
% matches Times optically. The journal keeps the same pairing.
\renewcommand{\rmdefault}{ptm}
\usepackage[scaled=0.92]{helvet}
\renewcommand{\ttdefault}{pcr}

\usepackage{cite}
\usepackage{amsmath,amssymb}
\usepackage{array}
\usepackage{booktabs}
\usepackage{algorithm}
\usepackage{algorithmic}
\usepackage{tikz}
\usepackage{pgfplots}
\pgfplotsset{compat=1.18}
\usetikzlibrary{arrows.meta,positioning,calc,fit,backgrounds}
\usepackage{microtype}
\usepackage{etoolbox}
\usepackage{enumitem}
\usepackage[font=small, labelsep=colon, skip=4pt]{caption}
\captionsetup[table]{position=top}

\makeatletter

% Section heads as in the survey. 12 pt bold, then 11 pt bold, numbered
% "1 Introduction" and "1.1 Title". A head that runs to two lines is set
% ragged right, so its first line is not stretched.
\renewcommand\section{\@startsection{section}{1}{\z@}%
  {-12pt \@plus -3pt \@minus -2pt}{5pt \@plus 1pt}%
  {\normalfont\large\bfseries\raggedright}}
\renewcommand\subsection{\@startsection{subsection}{2}{\z@}%
  {-9pt \@plus -2pt \@minus -1pt}{4pt \@plus 1pt}%
  {\normalfont\fontsize{11}{13}\selectfont\bfseries\raggedright}}

% Footnotes are 9 pt, as in both model papers. The patch must come before
% hyperref, which wraps \@footnotetext.
\patchcmd{\@footnotetext}{\footnotesize}{\small}{}%
  {\PackageWarning{preamble}{Footnote size patch failed}}

% The title block. The title sits between two 1 pt rules, as in KBQA-o1. The
% author names are 12 pt bold and the affiliation 12 pt roman, as in the
% survey. The notes go at the foot of the first column without a number.
\newcommand\affiliation[1]{\gdef\agr@affiliation{#1}}
\newcommand\agr@titlenotes{}
\newcommand\titlenote[1]{\g@addto@macro\agr@titlenotes{\footnotetext{#1}}}
\renewcommand\maketitle{%
  \twocolumn[{%
    \centering
    \hrule height 1pt
    \vskip 0.25in
    {\Large\bfseries\@title\par}%
    % KBQA-o1 puts the rule 16 pt below the last title baseline. The depth of
    % that line already gives 2.6 pt of it.
    \vskip 13.2pt
    \hrule height 1pt
    \vskip 0.3in
    {\large\bfseries\@author\par}%
    \vskip 4pt
    {\large\agr@affiliation\par}%
    \vskip 0.35in
  }]%
  \thispagestyle{plain}%
  \begingroup
    \renewcommand\thefootnote{}%
    \agr@titlenotes
  \endgroup}

% The abstract sits in the first column under a centred 12 pt bold head, set
% in from both sides by 20 pt, as in both model papers.
\renewenvironment{abstract}{%
  \par\noindent\hfil{\large\bfseries\abstractname}\hfil\par
  \vskip 6pt
  \list{}{\leftmargin 20pt \rightmargin 20pt \listparindent 1em
    \itemindent \z@ \parsep \z@ \topsep \z@ \partopsep \z@}%
  \item\relax}%
  {\endlist}
\newcommand\keywords[1]{%
  \begin{list}{}{\leftmargin 20pt \rightmargin 20pt \itemindent \z@
    \parsep \z@ \topsep 6pt \partopsep \z@}%
  \item\relax\textbf{Keywords:} #1%
  \end{list}\vskip 6pt}

% Reference entries are separated by a small gap rather than the list default.
\apptocmd{\thebibliography}{\setlength{\itemsep}{1.5pt \@plus 0.5pt}%
  \setlength{\parsep}{\z@}}{}%
  {\PackageWarning{preamble}{Bibliography spacing patch failed}}

% IEEEtran.bst reads its options from a control entry that \bstctlcite cites.
% IEEEtran.cls defines the command. This is its definition, for other classes.
\def\bstctlcite{\@ifnextchar[{\@bstctlcite}{\@bstctlcite[@auxout]}}
\def\@bstctlcite[#1]#2{\@bsphack
  \@for\@citeb:=#2\do{%
    \edef\@citeb{\expandafter\@firstofone\@citeb}%
    \if@filesw\immediate\write\csname #1\endcsname{\string\citation{\@citeb}}\fi}%
  \@esphack}

\makeatother

% Floats may fill more of a page than the article defaults allow, as in
% IEEEtran. Without this, most of the wide tables would go to float pages.
\renewcommand{\topfraction}{0.9}
\renewcommand{\dbltopfraction}{0.9}
\renewcommand{\textfraction}{0.1}
\renewcommand{\floatpagefraction}{0.8}
\renewcommand{\dblfloatpagefraction}{0.8}
\setcounter{topnumber}{3}
\setcounter{dbltopnumber}{2}
\setcounter{totalnumber}{4}

% A column that breaks early before a float or a head is left a line short
% rather than stretched. Stretching opened blank lines between paragraphs.
\raggedbottom

% Lists are set close to the text, as in the model papers.
\setlist{topsep=3pt plus 1pt, itemsep=1pt, parsep=0pt, partopsep=0pt}
\setlist[itemize,enumerate]{leftmargin=1.5em}

\usepackage[hidelinks]{hyperref}
\usepackage[capitalise,nameinlink]{cleveref}

% URLs in the text face, as in the survey.
\urlstyle{same}

% Two columns are narrow. A little emergency stretch lets TeX set a line with
% a long identifier in it rather than push it into the margin.
\setlength{\emergencystretch}{1.5em}

% Names split only where their parts join, or not at all.
\hyphenation{WebQSP Complex-Web-Questions GraphRAG HippoRAG StructGPT}

% Fixed-width table columns set ragged right, as in the book.
\newcolumntype{L}[1]{>{\raggedright\arraybackslash}p{#1}}

% Figure palette and plot style, copied from thesis_book/buetcsepgthesis.sty
% so that the generated figures render here exactly as they do in the book.
% Okabe-Ito colours, and every series also differs by mark shape, so a
% greyscale print keeps them apart.
\definecolor{agrNode}{HTML}{0072B2}
\definecolor{agrSup}{HTML}{009E73}
\definecolor{agrUns}{HTML}{D55E00}
\pgfplotsset{
  agrplot/.style={
    width=0.46\textwidth,
    height=52mm,
    tick label style={font=\scriptsize},
    label style={font=\small},
    title style={font=\small\bfseries},
    legend style={font=\scriptsize, draw=none, fill=none},
    grid=major,
    grid style={gray!25, very thin},
    axis line style={gray!60},
    tick align=outside,
    tick style={gray!60},
  }}

% cleveref's names. Figures and tables are named in full, as in the survey.
\crefname{algorithm}{Algorithm}{Algorithms}
\Crefname{algorithm}{Algorithm}{Algorithms}
\crefname{equation}{Eq.}{Eqs.}
\Crefname{equation}{Eq.}{Eqs.}

\endinput
